\chapter{Creio no perdão dos pecados}

O décimo artigo do Credo confessa a fé no \textbf{perdão dos pecados}, associado à fé no Espírito Santo, na Igreja e na comunhão dos santos. Cristo ressuscitado, ao dar o Espírito Santo aos apóstolos, conferiu-lhes o poder divino de perdoar os pecados: ``Recebei o Espírito Santo. Àqueles a quem perdoardes os pecados, ser-lhes-ão perdoados; àqueles a quem lhos retiverdes, ser-lhes-ão retidos'' (Jo 20,22-23). O Catecismo da Igreja Católica, nos parágrafos 976 a 987, explica que o Batismo é o primeiro e principal sacramento do perdão, mas a Penitência é necessária para os que caem depois do Batismo. Esta fé afirma a misericórdia infinita de Deus, que reconcilia os pecadores consigo e com a Igreja.

\section{Introdução (976)}

O Símbolo dos Apóstolos liga a fé no perdão dos pecados à fé no Espírito Santo, na Igreja e na comunhão dos santos. É pelo Espírito Santo que Cristo confiou aos apóstolos o poder de perdoar pecados. O Catecismo tratará explicitamente do perdão através do Batismo, da Penitência e dos outros sacramentos, especialmente da Eucaristia, mas aqui apresenta os factos essenciais.

\section{Um só Batismo para o perdão dos pecados (977--980)}

Nosso Senhor ligou o perdão dos pecados à fé e ao Batismo: ``Quem crer e for baptizado será salvo'' (Mc 16,16). O \textbf{Batismo é o primeiro e principal sacramento do perdão dos pecados}, porque une-nos a Cristo, que morreu pelos nossos pecados e ressuscitou para a nossa justificação, para que ``também nós possamos caminhar numa vida nova'' (Rm 6,4).

Ao professarmos a fé no Baptismo que nos purificou, o perdão recebido foi tão pleno e completo que nada restou a apagar: nem o pecado original, nem as ofensas cometidas pela nossa vontade, nem pena a sofrer para expiá-las. Contudo, a graça do Batismo não livra ninguém de toda a fraqueza da natureza. Pelo contrário, devemos ainda combater os movimentos da concupiscência que nos levam ao mal.

Nesta batalha contra a inclinação para o mal, quem poderia ser bastante forte e vigilante para escapar a toda ferida do pecado? Se a Igreja tem o poder de perdoar pecados, o Batismo não pode ser o único meio de usar as chaves do Reino dos Céus recebidas de Jesus Cristo. A Igreja deve perdoar a todos os penitentes as suas ofensas, mesmo que pequem até o último momento da vida.

\section{A reconciliação pelo Sacramento da Penitência e o poder da Igreja (981--987)}

É pelo \textbf{sacramento da Penitência} que os batizados se reconciliam com Deus e com a Igreja. Chamado pelos Santos Padres de ``Batismo laborioso'', este sacramento é necessário para a salvação dos que caem após o Batismo, assim como o Batismo o é para os que ainda não nasceram de novo.

Após a Ressurreição, Cristo enviou os apóstolos para que ``se pregasse em seu nome a penitência e o perdão dos pecados a todas as nações'' (Lc 24,47). Os apóstolos e os seus sucessores exercem este ``ministério da reconciliação'', anunciando o perdão de Deus, chamando à conversão e comunicando-o pelo Batismo e pela Penitência, pelo poder das chaves recebido de Cristo.

Não há ofensa, por grave que seja, que a Igreja não possa perdoar. ``Não há ninguém, por mais perverso e culpado que seja, que não possa esperar confiadamente o perdão, desde que a sua contrição seja sincera''. Cristo, que morreu por todos, quer que na sua Igreja as portas do perdão estejam sempre abertas a quem se arrepende.

O Catecismo esforça-se por despertar nos fiéis a fé na grandeza do dom de Cristo ressuscitado à sua Igreja: a missão e o poder de perdoar pecados pelo ministério dos apóstolos e seus sucessores. Os sacerdotes receberam de Deus um poder que nem aos anjos concedeu. Sem o perdão dos pecados na Igreja, não haveria esperança de vida eterna.

O Credo liga o ``perdão dos pecados'' à profissão de fé no Espírito Santo, pois o Cristo ressuscitado confiou aos apóstolos este poder ao dar-lhes o Espírito. O Batismo é o primeiro sacramento do perdão: une-nos a Cristo morto e ressuscitado e dá-nos o Espírito Santo.